%! Periodic Phenomena — Unit 6, Lesson 6.1 — Lesson Plan
\documentclass[10pt]{article}
\usepackage{precalculus-article}
\usepackage{precalculus-boxes}
\usepackage{graphicx}
\graphicspath{{images/}}
\newcommand{\TallMath}[1]{$\displaystyle #1\rule[-1.4em]{0pt}{3.2em}$}
\newcommand{\CourseName}{Precalculus}
\newcommand{\MeetingLength}{55 minutes}
\newcommand{\UnitNumberName}{Unit 6: Trigonometric Functions \quad}
\newcommand{\LessonNumberName}{Lesson 6.1: Periodic Phenomena}

\begin{document}

%----------------------------------------------------------------------
% Title Block
%----------------------------------------------------------------------
\begin{center}
  {\Large\bfseries \CourseName} \\[4pt]
  {\small\bfseries \UnitNumberName \LessonNumberName \quad (2 instructional periods, \MeetingLength\ each)}
\end{center}

%----------------------------------------------------------------------
% Primary Objective
%----------------------------------------------------------------------
\begin{tcolorbox}[colback=sky, colframe=navy, boxrule=0.9pt, arc=2mm,
    left=2mm, right=2mm, top=1.5mm, bottom=1.5mm]
  \textbf{Primary Objective:} Students will construct graphs of periodic
  relationships from verbal and tabular representations, and describe key
  characteristics — period, intervals of increase/decrease, concavity, and
  rate of change — explaining why each characteristic repeats in every period.
  \textit{(Big Idea: Functions)}
\end{tcolorbox}

\vspace{0.08in}

%----------------------------------------------------------------------
% Priority Ideas & Skills
%----------------------------------------------------------------------
\begin{skillbox}[Priority Ideas \& Skills]{goldbox}
\begin{tabularx}{\linewidth}{@{} p{0.46\linewidth} X @{}}
\textbf{AP Skills} & \textbf{Key Understandings (EKs)} \\
\midrule
\textbf{2.B} \textit{Construct equivalent representations} ---
  Build a complete graph of a periodic relationship by repeating one cycle
  from a verbal or tabular description.
&
\textbf{EK 3.1.A.1}: A periodic relationship exists when, as input increases,
output values demonstrate a repeating pattern over successive equal-length
intervals. \\[4pt]

\textbf{3.A} \textit{Describe characteristics} ---
  Identify and articulate the period, intervals of increase/decrease,
  concavity, and rate of change of a periodic function and explain why
  these repeat.
&
\textbf{EK 3.1.A.2}: The full graph of a periodic relationship can be
constructed from one cycle by repeating that cycle. \\[4pt]

&
\textbf{EK 3.1.B.1}: The period $k$ is the smallest positive value such
that $f(x+k)=f(x)$ for all $x$ in the domain; any interval of width $k$
completely determines the function. \\[4pt]

&
\textbf{EK 3.1.B.2}: Period can be estimated by finding where the output
pattern first begins to repeat. \\[4pt]

&
\textbf{EK 3.1.B.3}: Periodic functions share characteristics with other
functions (increase/decrease, concavity, rate of change), but in periodic
functions every characteristic found in one period appears in every period. \\
\end{tabularx}
\end{skillbox}

\vspace{0.08in}

%----------------------------------------------------------------------
% Vocabulary
%----------------------------------------------------------------------
\begin{skillbox}[{Vocabulary, Concepts, \& Theorems}]{greenbox}
\begin{tabularx}{\linewidth}{@{} p{0.24\linewidth} X @{}}
\textbf{Term} & \textbf{Definition / Note} \\
\midrule
periodic relationship &
  A relationship between two quantities in which, as input increases,
  output values repeat over successive equal-length input intervals. \\[4pt]
period &
  The smallest positive input length $k$ such that $f(x+k)=f(x)$ for all
  $x$ in the domain; the length of one complete cycle. \\[4pt]
cycle &
  One complete repetition of the repeating pattern of a periodic relationship. \\[4pt]
frequency (informal) &
  How often the pattern repeats per unit of input; informally the reciprocal
  of the period (not formally assessed in 3.1). \\
\end{tabularx}
\end{skillbox}

\vspace{0.08in}

%----------------------------------------------------------------------
% Activate Prior Knowledge
%----------------------------------------------------------------------
\begin{skillbox}[Activate Prior Knowledge \& Spiral Review]{sky}
\begin{tabularx}{\linewidth}{@{}X c@{}}
\begin{minipage}[t]{0.96\linewidth}\vspace{0pt}
\textbf{Spiral topics activated during Warm-Up (first 8 minutes):}
\begin{itemize}[leftmargin=1.4em, itemsep=2pt]
  \item \textbf{Function notation}: Evaluate $f(3)$ and $f(-1)$ from a
    provided table --- reinforces reading output values, which students
    use to spot repeating patterns.
  \item \textbf{Identifying patterns in tables}: Project a table of values
    and prompt students to describe what they notice. Connects to
    arithmetic/geometric sequences from Units 1--2; extends pattern
    recognition to repeating output values.
  \item \textbf{Average rate of change}: Review
    $\dfrac{f(b)-f(a)}{b-a}$ over a given interval --- students will
    apply this within and across periods in the activity.
\end{itemize}
\end{minipage}
&
\end{tabularx}
\end{skillbox}

\vspace{0.08in}

%----------------------------------------------------------------------
% Hook
%----------------------------------------------------------------------
\begin{skillbox}[Hook --- Entry Event]{sky}
\textbf{Scenario:} Display a table or projected image of 3 years of monthly
average electricity bills (or ocean tidal heights) — a dataset that goes up,
peaks, comes back down, and then repeats the same arc.

\textbf{Discussion prompts:}
\begin{enumerate}[leftmargin=1.4em, itemsep=3pt]
  \item \textit{``What pattern do you see in the data?''}
  \item \textit{``How is this pattern different from the linear or exponential
    functions we've studied? What doesn't fit those models?''}
  \item \textit{``If someone told you the month, could you predict the bill 5
    years from now? 10 years from now? What would you need to know?''}
\end{enumerate}

\textbf{Key tension to surface:} Unlike linear or exponential growth,
the values \textit{rise, fall, and repeat} on a fixed schedule. No constant
rate or fixed multiplier captures this. Students need a new idea --- a
\textit{repeating} function.

\textbf{Transition:} ``Today we name this a \textbf{periodic relationship}.
By the end of class you will be able to identify one, estimate how long a
single cycle lasts, and construct a complete graph from just one cycle.''
\end{skillbox}

\vspace{0.08in}

%----------------------------------------------------------------------
% Lesson
%----------------------------------------------------------------------
\begin{skillbox}[Lesson --- Instruction Sequence]{sky}
\begin{multicols}{2}

\textbf{Step 1 --- Identify Periodic Relationships} (10 min)

Present a table of average monthly high temperatures for a city over
24 months. Ask students to observe and describe the output pattern.
Formally define ``periodic relationship'' (EK 3.1.A.1): the pattern must
\emph{repeat exactly} over \emph{equal-length} input intervals.
Contrast with non-periodic up-and-down patterns (e.g., irregular stock prices)
to sharpen the definition.

\columnbreak

\textbf{Step 2 --- Estimating the Period} (10 min)

Model how to estimate the period from the temperature table: find the
smallest input length $k$ such that the output pattern restarts.
Formalize: $f(x+k)=f(x)$ (EK 3.1.B.1--B.2).
Partner practice: present a second table (height of a point on a spinning
wheel sampled every second) and have pairs estimate the period and
justify their answer.

\end{multicols}

\begin{multicols}{2}

\textbf{Step 3 --- Construct Graph from One Cycle} (12 min)

Model: plot one cycle of the temperature data on the coordinate grid in the
guided notes. Then \emph{extend} the graph by repeating the cycle exactly ---
same shape, same width (EK 3.1.A.2). Students practice on a pre-printed
grid in their guided notes: one cycle is partially filled, students
extend it for two additional periods.

\columnbreak

\textbf{Step 4 --- Key Characteristics} (10 min)

Return to the constructed graph. Identify and label:
\begin{itemize}[leftmargin=1em, itemsep=1pt]
  \item intervals of increase and decrease
  \item concavity (concave up / concave down regions)
  \item approximate rate of change on a given interval
\end{itemize}
Emphasize EK 3.1.B.3: every characteristic from one period appears in
every period. Close: \textit{``Why does knowing one period completely
determine the whole function?''}

\end{multicols}

\smallskip
\textbf{Pacing note:} Steps 1--2 in Period 1 (days 1);
Steps 3--4 + activity in Period 2 (day 2).
\end{skillbox}

\vspace{0.08in}

%----------------------------------------------------------------------
% Active Monitoring
%----------------------------------------------------------------------
\begin{skillbox}[Active Monitoring]{redbox}
\begin{itemize}[leftmargin=1.4em, itemsep=3pt]
  \item \textbf{Period vs.\ amplitude confusion}: Students often report the
    range of outputs (how high/low) instead of the input length of one cycle.
    Prompt: \textit{``The period measures \underline{input} --- how far along
    the $x$-axis before it repeats. What input value marks the end of one cycle?''}
  \item \textbf{Pattern vs.\ periodic}: Any up-and-down table may look periodic.
    Ask: \textit{``Is the interval between repeats always the same length?
    Is the shape identical in every cycle?''}
  \item \textbf{Graph construction errors}: Students stretch or compress
    subsequent cycles. Remind: each cycle is \emph{identical} in width and
    shape to the first.
  \item \textbf{Characteristics}: Verify students label increase/decrease with
    input intervals (not output values) and articulate \emph{why} characteristics
    repeat (not just that they do).
\end{itemize}
\end{skillbox}

\vspace{0.08in}

%----------------------------------------------------------------------
% Group Work & Differentiation
%----------------------------------------------------------------------
\begin{skillbox}[Group Work \& Differentiation]{redbox}
Students work in groups of 3--4 on the tiered activity (assign by teacher
or allow student self-selection with guidance).

\begin{itemize}[leftmargin=1.4em, itemsep=4pt]
  \item \textbf{Tier R (Remediate)}: Given a completed table with the
    repeating block highlighted, students identify whether a relationship
    is periodic, name the period from the table, and describe behavior in
    one cycle using sentence frames provided on the worksheet.
  \item \textbf{Tier A (Accelerate)}: Given a verbal description of a periodic
    phenomenon (e.g., a water wheel with given heights at each quarter-turn),
    students build a table of values for two full periods, estimate the period,
    and describe intervals of increase and decrease.
  \item \textbf{Tier E (Extend)}: Students compare two periodic phenomena
    provided as tables. They determine if each is periodic, identify both
    periods, and compute average rates of change on specified intervals to
    compare the two functions quantitatively.
\end{itemize}
\end{skillbox}

\vspace{0.08in}

%----------------------------------------------------------------------
% Individual Work & Assessment
%----------------------------------------------------------------------
\begin{skillbox}[Individual Work \& Assessment]{redbox}
Exit ticket administered independently (final 5 minutes of Period 2).

\begin{enumerate}[leftmargin=1.4em, itemsep=4pt]
  \item A short 8-value table is provided. Students determine whether the
    relationship is periodic and estimate the period with justification.
  \item Open response: \textit{``Name one characteristic of a periodic function
    that repeats in every period and explain in one sentence why this must
    be so.''}
\end{enumerate}

\textbf{Data use:} Sort exit tickets into three piles --- correctly identified
period, incorrect period, non-attempt. Use groupings to inform the opening
of Lesson 3.2 (brief re-teach or accelerated start).
\end{skillbox}

\vspace{0.08in}

%----------------------------------------------------------------------
% Reinforcement & Extension
%----------------------------------------------------------------------
\begin{skillbox}[Reinforcement \& Extension]{goldbox}
\textbf{Homework (Lesson 6.1):} 5--6 problems covering: identifying periodic
relationships from tables and descriptions, estimating periods, describing
characteristics within a period, and 2 AP-style multiple-choice items.
An optional extension challenge asks students to compare average rates of
change on two different intervals within a single period.

\textbf{Preview of Lesson 3.2 --- Sine, Cosine, and Tangent:}
Ask students to consider: \textit{``What if a periodic relationship could be
described with a \emph{formula} rather than a table? How would that be more
powerful?''} Lesson 3.2 introduces sine and cosine as the canonical models
for periodic phenomena and directly connects today's period, increase/decrease,
and concavity language to amplitude, period, and midline of sinusoidal functions.
\end{skillbox}

\end{document}
